\documentclass[10pt]{article}

\usepackage[margin=1in]{geometry}
\usepackage{amsmath,amssymb,amsfonts,mathtools,bm}
\usepackage{booktabs,multirow,array}
\usepackage{algorithm}
\usepackage{algpseudocode}
\usepackage{xcolor}
\usepackage{hyperref}
\usepackage{microtype}
\usepackage{enumitem}
\usepackage{graphicx}
\usepackage{caption}
\usepackage{cleveref}

\hypersetup{colorlinks=true,linkcolor=blue,citecolor=blue,urlcolor=blue}
\newcommand{\E}{\mathbb{E}}
\newcommand{\KL}{D_{\mathrm{KL}}}
\newcommand{\TV}{D_{\mathrm{TV}}}
\newcommand{\V}{\mathcal{V}}
\newcommand{\DF}{\mathcal{D}_{F}}
\newcommand{\DR}{\mathcal{D}_{R}}
\newcommand{\M}{\mathcal{M}}
\newcommand{\sg}{\operatorname{sg}}
\newcommand{\softmax}{\operatorname{softmax}}
\newcommand{\1}{\mathbb{I}}

\title{TIDE-U: Trajectory-Consistent Importance-Weighted Distributional Erasure\\for Masked Diffusion Language Models}
\author{Research Proposal}
\date{}

\begin{document}
\maketitle

\begin{abstract}
We propose TIDE-U, a loss-based unlearning method for masked diffusion language models (MDLMs). The method separates two goals: removing the probability of generating target knowledge and preserving the original model's fluent behavior on non-target outputs. Instead of only decreasing the likelihood of forget tokens, TIDE-U constructs a normalized replacement distribution that moves probability mass toward fluent, non-leaking alternatives supported by the original model. The replacement distribution is obtained by exponentially tilting the original reverse-denoising trajectory distribution with a semantic leakage cost. A Doob-transform derivation converts this sequence-level objective into token-level targets at individual masked positions and denoising timesteps. The resulting forward-KL loss has an explicit logit gradient: tokens with above-average downstream leakage risk are decreased, while safer alternatives are increased. We provide implementation-ready notation, objectives, algorithms, and practical approximations.
\end{abstract}

\section{Problem Setup}

Let an MDLM generate a length-$L$ sequence $x_0=(x_{0,1},\ldots,x_{0,L})\in\V^L$ from a prompt $c$. The reverse process starts from a masked sequence $x_{t_K}$ and iteratively denoises it:
\begin{equation}
\tau=(x_{t_K},x_{t_{K-1}},\ldots,x_{t_0}),\qquad x_{t_0}=x_0.
\end{equation}
The model trajectory distribution is
\begin{equation}
\Pi_{\theta}(\tau\mid c,\sigma)
=
\prod_{k=K}^{1}
P_{\theta,k}(x_{t_{k-1}}\mid x_{t_k},c,\sigma),
\label{eq:trajectory}
\end{equation}
where $\sigma$ is an unmasking or remasking schedule. The original model is denoted by $\theta_0$, and the unlearned model by $\theta$.

We assume two datasets:
\begin{itemize}[leftmargin=*]
    \item Forget data $\DF$: prompts and target knowledge that should be removed.
    \item Retain data $\DR$: ordinary prompts on which behavior should remain close to the original model.
\end{itemize}

A semantic leakage function
\begin{equation}
C_F(x_0,c)\in[0,1]
\end{equation}
measures how strongly the generated response reveals the forgotten knowledge. In the simplest case, $C_F\in\{0,1\}$ is a binary semantic verifier. It should detect paraphrases, aliases, reordered facts, and indirect statements, not only exact token matches.

\subsection{Two separate objectives}

\paragraph{Forgetting.}
The model should have low probability of producing the target knowledge:
\begin{equation}
\E_{c\sim\DF}
\E_{\tau\sim\Pi_\theta(\cdot\mid c)}
[C_F(x_0,c)]
\leq \varepsilon.
\label{eq:forget_constraint}
\end{equation}

\paragraph{Preservation.}
Among models satisfying \cref{eq:forget_constraint}, the reverse process should remain as close as possible to the original model on non-target behavior. On retain prompts, we use token-level forward KL:
\begin{equation}
\mathcal{L}_{\mathrm{retain}}
=
\E_{\substack{c_r\sim\DR,\,t,\,x_t}}
\left[
\rho(t)\frac{1}{|\M_t|}
\sum_{i\in\M_t}
\KL\!\left(
 p_{\theta_0,i}(\cdot\mid c_r,x_t)
 \Vert
 p_{\theta,i}(\cdot\mid c_r,x_t)
\right)
\right],
\label{eq:retain_loss}
\end{equation}
where $\M_t$ is the set of masked positions and $\rho(t)$ is the standard MDLM timestep weight.

\section{Notation}

\begin{table}[t]
\centering
\caption{Main notation used by TIDE-U.}
\label{tab:notation}
\small
\begin{tabular}{p{0.20\linewidth}p{0.72\linewidth}}
\toprule
Symbol & Meaning \\
\midrule
$c$ & Input prompt. \\
$x_0$ & Final fully denoised response. \\
$x_t$ & Partially masked response at diffusion time $t$. \\
$\tau$ & Complete reverse-denoising trajectory. \\
$\sigma$ & Unmasking/remasking schedule or reveal order. \\
$\theta_0$ & Frozen original MDLM parameters. \\
$\theta$ & Trainable unlearned MDLM parameters. \\
$p_{\theta,i}(v\mid c,x_t)$ & Token probability for vocabulary item $v$ at masked position $i$. \\
$\ell_{i,t,v}$ & Logit for token $v$ at position $i$ and time $t$. \\
$C_F(x_0,c)$ & Semantic leakage score for the forgotten knowledge. \\
$\beta$ & Forgetting strength / Lagrange multiplier. \\
$h_{i,t}(v\mid s)$ & Desirability of choosing token $v$ and then completing safely. \\
$r_{i,t}(v\mid s)$ & Probability that choosing $v$ eventually produces target knowledge. \\
$q^\star_{i,t}(v\mid s)$ & Safe replacement token distribution. \\
$\Delta_{i,t}$ & Token knowledge-attribution score. \\
$G_{i,t}$ & Required redistribution magnitude. \\
$g_{i,t}$ & Soft or hard importance gate. \\
$\rho(t)$ & MDLM timestep weighting factor. \\
$G$ & Number of rollout completions per candidate token. \\
$\varepsilon$ & Desired maximum semantic leakage probability. \\
\bottomrule
\end{tabular}
\end{table}

\section{Minimum-Change Forgetting Distribution}

For a forget prompt $c$, define the regularized objective over trajectory distributions $\Pi$:
\begin{equation}
J_\beta(\Pi)
=
\E_{\tau\sim\Pi}[C_F(x_0,c)]
+
\frac{1}{\beta}\KL(\Pi\Vert\Pi_0).
\label{eq:regularized}
\end{equation}
The first term removes target knowledge. The second term preserves the original reverse process.

Define
\begin{equation}
Z_\beta(c)
=
\E_{\tau\sim\Pi_0(\cdot\mid c)}
\left[e^{-\beta C_F(x_0,c)}\right]
\end{equation}
and
\begin{equation}
\Pi_\beta^\star(\tau\mid c)
=
\frac{
\Pi_0(\tau\mid c)e^{-\beta C_F(x_0,c)}
}{Z_\beta(c)}.
\label{eq:tilted_trajectory}
\end{equation}
Then
\begin{align}
\KL(\Pi\Vert\Pi_\beta^\star)
&=
\E_{\Pi}\left[
\log\frac{\Pi(\tau)}{\Pi_0(\tau)e^{-\beta C_F(\tau)}/Z_\beta}
\right]\\
&=
\KL(\Pi\Vert\Pi_0)
+
\beta\E_{\Pi}[C_F]
+
\log Z_\beta.
\end{align}
Therefore,
\begin{equation}
J_\beta(\Pi)
=
\frac{1}{\beta}\KL(\Pi\Vert\Pi_\beta^\star)
-
\frac{1}{\beta}\log Z_\beta,
\end{equation}
and the unique optimum is
\begin{equation}
\arg\min_\Pi J_\beta(\Pi)=\Pi_\beta^\star.
\end{equation}

This gives a simple interpretation: keep the original trajectory probabilities, but multiply leaking trajectories by a smaller weight.

\subsection{Exact leakage reduction for a binary verifier}

Let
\begin{equation}
m_0(c)=\Pr_{x_0\sim p_{\theta_0}(\cdot\mid c)}[C_F(x_0,c)=1].
\end{equation}
For binary $C_F$,
\begin{equation}
Z_\beta(c)=1-m_0(c)+e^{-\beta}m_0(c),
\end{equation}
and the target leakage probability is
\begin{equation}
m_\beta^\star(c)
=
\frac{e^{-\beta}m_0(c)}{1-m_0(c)+e^{-\beta}m_0(c)}.
\label{eq:ideal_leakage}
\end{equation}
To guarantee $m_\beta^\star(c)\leq\varepsilon$, choose
\begin{equation}
\beta
\geq
\log\frac{m_0(c)(1-\varepsilon)}{\varepsilon(1-m_0(c))}.
\end{equation}
In practice, $\beta$ can be updated by dual ascent using the measured leakage rate.

\subsection{Why the target remains fluent}

The final output distribution is
\begin{equation}
p_\beta^\star(x_0\mid c)
=
\frac{p_{\theta_0}(x_0\mid c)e^{-\beta C_F(x_0,c)}}{Z_\beta(c)}.
\end{equation}
For any two non-leaking outputs $y_1,y_2$,
\begin{equation}
\frac{p_\beta^\star(y_1\mid c)}{p_\beta^\star(y_2\mid c)}
=
\frac{p_{\theta_0}(y_1\mid c)}{p_{\theta_0}(y_2\mid c)}.
\label{eq:nonleak_odds}
\end{equation}
Thus, among safe outputs, the original ranking is unchanged. The method does not reward gibberish or fixed refusals; it preserves the original model's preference for natural alternatives.

\section{MDLM Token-Level Safe Target}

Let the current state be
\begin{equation}
s=(c,x_t,t).
\end{equation}
At a masked position $i$, consider choosing token $v$. Define the desirability
\begin{equation}
h_{i,t}(v\mid s)
=
\E_{\Pi_0}
\left[
 e^{-\beta C_F(X_0,c)}
 \mid s, X_i=v
\right].
\label{eq:desirability}
\end{equation}
The normalized safe token target is
\begin{equation}
q_{i,t}^\star(v\mid s)
=
\frac{
 p_{\theta_0,i}(v\mid s)h_{i,t}(v\mid s)
}{
 \sum_{u\in\V}p_{\theta_0,i}(u\mid s)h_{i,t}(u\mid s)
}.
\label{eq:safe_target}
\end{equation}

For binary leakage, define
\begin{equation}
r_{i,t}(v\mid s)
=
\Pr_{\Pi_0}
\left[
C_F(X_0,c)=1\mid s,X_i=v
\right].
\end{equation}
Then
\begin{align}
h_{i,t}(v\mid s)
&=(1-r_{i,t}(v\mid s))+e^{-\beta}r_{i,t}(v\mid s)\\
&=1-(1-e^{-\beta})r_{i,t}(v\mid s).
\label{eq:h_binary}
\end{align}
A token is therefore downweighted only when choosing it at this exact state and timestep makes eventual semantic leakage likely.

\section{Estimating the Safe Target with Rollouts}

Exact computation of \cref{eq:desirability} is intractable. We estimate it by completing the denoising process from a candidate token.

For each selected candidate $v$, sample
\begin{equation}
X_0^{(g)}\sim\Pi_0(\cdot\mid s,X_i=v),
\qquad g=1,\ldots,G.
\end{equation}
Use
\begin{equation}
\widehat h_{i,t}(v\mid s)
=
\frac{1}{G}\sum_{g=1}^{G}
\exp\left[-\beta C_F(X_0^{(g)},c)\right].
\label{eq:h_estimator}
\end{equation}
Then construct
\begin{equation}
\widehat q_{i,t}^\star(v\mid s)
=
\frac{
 p_{\theta_0,i}(v\mid s)\widehat h_{i,t}(v\mid s)
}{
 \sum_{u\in\mathcal C_{i,t}}
 p_{\theta_0,i}(u\mid s)\widehat h_{i,t}(u\mid s)
},
\label{eq:qhat}
\end{equation}
where $\mathcal C_{i,t}$ is a small candidate set. For easy implementation, use the union of:
\begin{itemize}[leftmargin=*]
    \item the top-$K_0$ tokens from the original model;
    \item the top-$K_\theta$ tokens from the current model;
    \item known forget-target tokens and aliases;
    \item the ground-truth token, if available.
\end{itemize}
All vocabulary mass outside $\mathcal C_{i,t}$ can be grouped into a residual bucket or left unchanged through the preservation loss.

\section{Selecting Important Tokens and Timesteps}

We avoid modifying every position. For a leaked or ground-truth response $y$, define a prompt-conditioned attribution score
\begin{equation}
\Delta_{i,t}
=
\left|
\log p_{\theta_0,i}(y_i\mid c,x_t)
-
\log p_{\theta_0,i}(y_i\mid \widetilde c,x_t)
\right|,
\label{eq:attribution}
\end{equation}
where $\widetilde c$ masks or removes the target entity/relation while preserving the rest of the prompt.

Define state entropy
\begin{equation}
H_{i,t}
=-\sum_{v\in\V}p_{\theta_0,i}(v\mid s)
\log p_{\theta_0,i}(v\mid s).
\end{equation}
After normalization within the sequence,
\begin{equation}
\phi_{i,t}
=\alpha\overline\Delta_{i,t}+(1-\alpha)\overline H_{i,t}.
\end{equation}

The amount of distributional change requested by the safe target is
\begin{equation}
G_{i,t}
=
\left\|q_{i,t}^\star-p_{\theta_0,i}\right\|_1.
\label{eq:redistribution}
\end{equation}
The final importance score is
\begin{equation}
S_{i,t}=\phi_{i,t}G_{i,t}.
\end{equation}
Use either a hard top-$r\%$ gate or a soft gate:
\begin{equation}
g_{i,t}=\sigma\!\left(\frac{S_{i,t}-\tau_g}{T_g}\right),
\label{eq:gate}
\end{equation}
where $\tau_g$ is a threshold and $T_g$ is a temperature.

\section{Training Objective}

\subsection{Safe redistribution loss}

At important forget positions, match the safe replacement distribution:
\begin{equation}
\mathcal L_F
=
\E_{\substack{c\sim\DF,\,\sigma,\,\tau,\,t}}
\left[
\rho(t)
\frac{1}{|\M_t|}
\sum_{i\in\M_t}
 g_{i,t}
\KL\!\left(
\sg[q_{i,t}^\star]\Vert p_{\theta,i}
\right)
\right].
\label{eq:forget_loss}
\end{equation}

\subsection{Context preservation on forget prompts}

At noncritical positions, preserve the original token distribution:
\begin{equation}
\mathcal L_{\mathrm{context}}
=
\E_{\DF,\sigma,\tau,t}
\left[
\rho(t)
\frac{1}{|\M_t|}
\sum_{i\in\M_t}
(1-g_{i,t})
\KL\!\left(
 p_{\theta_0,i}\Vert p_{\theta,i}
\right)
\right].
\label{eq:context_loss}
\end{equation}

\subsection{Total objective}

\begin{equation}
\boxed{
\mathcal L_{\mathrm{TIDE-U}}
=
\mathcal L_F
+\lambda_C\mathcal L_{\mathrm{context}}
+\lambda_R\mathcal L_{\mathrm{retain}}.
}
\label{eq:total_loss}
\end{equation}

\section{Explicit Logit-Gradient Effect}

Let
\begin{equation}
p_\theta(v)=\frac{e^{\ell_v}}{\sum_u e^{\ell_u}}
\end{equation}
for one selected state and position. For
\begin{equation}
\ell_{\mathrm{safe}}=\KL(q^\star\Vert p_\theta),
\end{equation}
we have
\begin{equation}
\frac{\partial \ell_{\mathrm{safe}}}{\partial\ell_v}
=p_\theta(v)-q^\star(v).
\label{eq:basic_gradient}
\end{equation}
At initialization, $p_\theta=p_{\theta_0}=p_0$. Let
\begin{equation}
\overline h=\sum_u p_0(u)h(u).
\end{equation}
Using \cref{eq:safe_target},
\begin{equation}
\left.
\frac{\partial\mathcal L_F}{\partial\ell_v}
\right|_{\theta=\theta_0}
=
\rho(t)g_{i,t}p_0(v)
\left(1-\frac{h(v)}{\overline h}\right).
\label{eq:target_gradient}
\end{equation}
Hence:
\begin{itemize}[leftmargin=*]
    \item If $h(v)<\overline h$, token $v$ is riskier than average. The gradient is positive, so gradient descent decreases its logit.
    \item If $h(v)>\overline h$, token $v$ is safer than average. The gradient is negative, so gradient descent increases its logit.
\end{itemize}

For binary leakage, with $\kappa=1-e^{-\beta}$ and $\overline r=\sum_u p_0(u)r(u)$,
\begin{equation}
h(v)=1-\kappa r(v),
\end{equation}
so
\begin{equation}
r(v)>\overline r
\Longrightarrow
\frac{\partial\mathcal L_F}{\partial\ell_v}>0.
\end{equation}
Thus the loss acts specifically on tokens whose choice at the current denoising state increases eventual leakage.

For a risky token $f$ and safe token $a$,
\begin{equation}
\log\frac{q^\star(f)}{q^\star(a)}
=
\log\frac{p_0(f)}{p_0(a)}
+
\log\frac{1-\kappa r(f)}{1-\kappa r(a)}.
\label{eq:odds}
\end{equation}
If $r(f)=1$ and $r(a)=0$,
\begin{equation}
\log\frac{q^\star(f)}{q^\star(a)}
=
\log\frac{p_0(f)}{p_0(a)}-\beta.
\end{equation}
The forget-token log odds are reduced by exactly $\beta$ relative to a safe alternative.

Also,
\begin{equation}
\sum_v\frac{\partial\mathcal L_F}{\partial\ell_v}=0,
\end{equation}
so probability removed from risky tokens is explicitly transferred to other tokens rather than destroyed.

\section{Why Forgetting Persists Across Denoising Steps}

Define a trajectory-level desirability
\begin{equation}
h_k(x_{t_k},c)
=
\E_{\Pi_0}
\left[
 e^{-\beta C_F(x_0,c)}
 \mid x_{t_k},c
\right].
\end{equation}
The target reverse transition is
\begin{equation}
P_{\beta,k}^\star(x_{t_{k-1}}\mid x_{t_k},c)
=
P_{0,k}(x_{t_{k-1}}\mid x_{t_k},c)
\frac{h_{k-1}(x_{t_{k-1}},c)}{h_k(x_{t_k},c)}.
\label{eq:doob}
\end{equation}
All timesteps use the same terminal semantic cost. Therefore, an early token is penalized if it creates a state from which the forgotten fact can be reconstructed later, and a late token is penalized if it completes the fact.

If the learned process satisfies
\begin{equation}
\delta_F
=
\KL(\Pi_\beta^\star\Vert\Pi_\theta),
\end{equation}
then the KL chain rule gives
\begin{equation}
\delta_F
=
\sum_{k=K}^{1}
\E_{x_{t_k}\sim\Pi_\beta^\star}
\left[
\KL\!\left(
P_{\beta,k}^\star(\cdot\mid x_{t_k},c)
\Vert
P_{\theta,k}(\cdot\mid x_{t_k},c)
\right)
\right].
\end{equation}
By Pinsker's inequality,
\begin{equation}
\Pr_{\Pi_\theta}(C_F=1)
\leq
m_\beta^\star(c)+\sqrt{\frac{\delta_F}{2}}.
\label{eq:bound}
\end{equation}
Therefore, fitting local transition targets across states encountered during inference controls the final semantic leakage probability.

\section{MDLM-Specific Implementation Details}

\subsection{Random reveal orders}

An MDLM can reveal tokens in many orders. Train with multiple random schedules $\sigma$ per forget example. A token's importance and safe target must depend on $(x_t,t)$, not only on its absolute position.

\subsection{Parallel unmasking}

For the simplest implementation and cleanest derivation, use one-token microsteps during target construction. The production sampler may still reveal multiple tokens in parallel. If training directly with a block $A$, the exact target is
\begin{equation}
q_A^\star(v_A\mid s)
\propto
P_{0,A}(v_A\mid s)h(s,v_A),
\end{equation}
which may contain dependencies among tokens. Independent token targets are an approximation when multiple positions are revealed simultaneously.

\subsection{Remasking}

If inference uses confidence-based remasking, generate training states with the same remasking rule. Otherwise, the training-state distribution may not match inference trajectories.

\section{Implementation-Ready Algorithms}

\begin{algorithm}[t]
\caption{Estimate a safe token target at one MDLM state}
\label{alg:safe_target}
\begin{algorithmic}[1]
\Require Frozen MDLM $\theta_0$, current MDLM $\theta$, prompt $c$, state $x_t$, position $i$, forgetting strength $\beta$, rollout count $G$, top-$K$ sizes $K_0,K_\theta$, semantic verifier $C_F$
\State $p_0\gets p_{\theta_0,i}(\cdot\mid c,x_t)$
\State $p_{\mathrm{cur}}\gets p_{\theta,i}(\cdot\mid c,x_t)$
\State $\mathcal C\gets \mathrm{TopK}(p_0,K_0)\cup\mathrm{TopK}(p_{\mathrm{cur}},K_\theta)$
\State Add known target tokens, aliases, and the reference token to $\mathcal C$
\For{each $v\in\mathcal C$}
    \State Create state $x_t^{(v)}$ by fixing position $i$ to token $v$
    \State $H_v\gets 0$
    \For{$g=1$ to $G$}
        \State Sample completion $X_0^{(g)}\sim\Pi_0(\cdot\mid c,x_t^{(v)})$
        \State $H_v\gets H_v+\exp[-\beta C_F(X_0^{(g)},c)]$
    \EndFor
    \State $\widehat h(v)\gets H_v/G$
    \State $\widetilde q(v)\gets p_0(v)\widehat h(v)$
\EndFor
\State $q^\star(v)\gets \widetilde q(v)/\sum_{u\in\mathcal C}\widetilde q(u)$ for $v\in\mathcal C$
\State Preserve or aggregate probability mass outside $\mathcal C$
\State \Return $q^\star$
\end{algorithmic}
\end{algorithm}

\begin{algorithm}[t]
\caption{TIDE-U training}
\label{alg:training}
\begin{algorithmic}[1]
\Require Original MDLM $\theta_0$, trainable copy $\theta$, forget data $\DF$, retain data $\DR$, verifier $C_F$, hyperparameters $\lambda_C,\lambda_R,\alpha,\tau_g,T_g,G$
\For{each training step}
    \State Sample forget minibatch $B_F\sim\DF$ and retain minibatch $B_R\sim\DR$
    \State Initialize $L_F\gets0$, $L_C\gets0$, $L_R\gets0$
    \For{each forget example $c$ in $B_F$}
        \State Sample a reverse schedule $\sigma$ and one or more timesteps $t$
        \State Generate or retrieve an inference-like state $x_t$
        \State Compute original and current token distributions at all $i\in\M_t$
        \State Compute attribution $\Delta_{i,t}$ using $c$ and counterfactual prompt $\widetilde c$
        \State Compute entropy $H_{i,t}$ and preliminary score $\phi_{i,t}$
        \State Select a small set of candidate positions using top attribution scores
        \For{each selected position $i$}
            \State $q_{i,t}^\star\gets$ \Call{EstimateSafeTarget}{$\theta_0,\theta,c,x_t,i,\beta,G$}
            \State $G_{i,t}\gets\|q_{i,t}^\star-p_{\theta_0,i}\|_1$
            \State $S_{i,t}\gets\phi_{i,t}G_{i,t}$
            \State $g_{i,t}\gets\sigma((S_{i,t}-\tau_g)/T_g)$
            \State $L_F\gets L_F+\rho(t)g_{i,t}\KL(\sg[q_{i,t}^\star]\Vert p_{\theta,i})$
            \State $L_C\gets L_C+\rho(t)(1-g_{i,t})\KL(p_{\theta_0,i}\Vert p_{\theta,i})$
        \EndFor
        \State Apply context preservation to unselected masked positions
    \EndFor
    \For{each retain prompt $c_r$ in $B_R$}
        \State Sample $t$ and construct masked state $x_t$
        \State $L_R\gets L_R+\rho(t)|\M_t|^{-1}\sum_{i\in\M_t}\KL(p_{\theta_0,i}\Vert p_{\theta,i})$
    \EndFor
    \State $L\gets L_F+\lambda_C L_C+\lambda_R L_R$
    \State Update $\theta\gets\theta-\eta\nabla_\theta L$
    \State Periodically estimate leakage $\widehat m$ and update $\beta\gets[\beta+\eta_\beta(\widehat m-\varepsilon)]_+$
\EndFor
\State \Return $\theta$
\end{algorithmic}
\end{algorithm}

\begin{algorithm}[t]
\caption{Low-cost practical variant}
\label{alg:cheap}
\begin{algorithmic}[1]
\Require Same inputs as \cref{alg:training}
\State Precompute target aliases and paraphrases for each forget concept
\State Use only the top $N_{\mathrm{pos}}$ positions by $\Delta_{i,t}$
\State Use $K_0=8$, $K_\theta=8$, and $G\in\{2,4\}$ rollouts
\State Cache $\widehat h_{i,t}(v)$ for similar states and refresh every $R$ updates
\State Use one random timestep and one random reveal order per example per step
\State Use LoRA or another parameter-efficient adapter for $\theta$
\State Train with \cref{eq:total_loss}; periodically regenerate current-model leakage examples
\end{algorithmic}
\end{algorithm}

\section{Recommended Starting Hyperparameters}

\begin{table}[t]
\centering
\caption{Practical starting values. These are initial defaults, not theoretically fixed constants.}
\label{tab:hparams}
\small
\begin{tabular}{lll}
\toprule
Hyperparameter & Starting value & Comment \\
\midrule
$N_{\mathrm{pos}}$ & 2--8 positions & Most important masked positions per state. \\
$K_0$ & 8--16 & Top original-model token candidates. \\
$K_\theta$ & 8--16 & Top current-model token candidates. \\
$G$ & 2--8 & More rollouts improve risk estimates. \\
$\beta$ & 2--6 & Update by dual ascent if possible. \\
$\lambda_C$ & 0.5--2.0 & Forget-prompt context preservation. \\
$\lambda_R$ & 0.5--4.0 & Retain behavior preservation. \\
$\alpha$ & 0.7 & More weight on prompt-conditioned attribution. \\
$\tau_g$ & median score & Stable initial gate threshold. \\
$T_g$ & 0.05--0.2 & Smaller values approach hard selection. \\
LoRA rank & 8--32 & Reduces collateral parameter changes. \\
Leakage target $\varepsilon$ & task-dependent & E.g., $0.01$ or $0.05$. \\
\bottomrule
\end{tabular}
\end{table}

\section{Minimal PyTorch-Style Loss}

The core loss can be implemented as follows:
\begin{verbatim}
# logits_cur: [B, L, V]
# probs_base: [B, L, V], detached
# q_safe:    [B, L, V], detached
# gate:      [B, L], in [0, 1]
# mask:      [B, L], 1 for masked prediction positions

logp_cur = logits_cur.log_softmax(dim=-1)

kl_safe = (q_safe * (q_safe.clamp_min(1e-8).log() - logp_cur)).sum(-1)
kl_base = (probs_base *
          (probs_base.clamp_min(1e-8).log() - logp_cur)).sum(-1)

forget_loss = ((mask * gate) * kl_safe).sum() / (mask * gate).sum().clamp_min(1)
context_loss = ((mask * (1-gate)) * kl_base).sum() / \
               (mask * (1-gate)).sum().clamp_min(1)

loss = forget_loss + lambda_context * context_loss + lambda_retain * retain_loss
\end{verbatim}

For a restricted candidate set, create $q_{\mathrm{safe}}$ only on candidate tokens and preserve all other token probabilities from the base model before renormalizing:
\begin{equation}
\widetilde q(v)=
\begin{cases}
p_0(v)\widehat h(v), & v\in\mathcal C,\\
p_0(v), & v\notin\mathcal C,
\end{cases}
\qquad
q^\star(v)=\frac{\widetilde q(v)}{\sum_u\widetilde q(u)}.
\label{eq:full_vocab_approx}
\end{equation}
This approximation is simple, normalized, and guarantees that non-candidate tokens are not arbitrarily changed.

\section{Evaluation}

A successful method must satisfy all three criteria:
\begin{enumerate}[leftmargin=*]
    \item Low semantic leakage under direct, paraphrased, indirect, multi-hop, and adversarial prompts.
    \item Low divergence from the original model on retain prompts and noncritical positions.
    \item Fluent, coherent outputs without a large increase in gibberish, truncation, repetition, or generic refusal rates.
\end{enumerate}

Recommended measurements include:
\begin{itemize}[leftmargin=*]
    \item semantic leakage probability under multiple reveal orders and samplers;
    \item probability mass on target tokens, semantic neighbors, safe alternatives, and the vocabulary tail;
    \item token KL to the original model on retain data;
    \item independent-LM perplexity and human coherence judgments;
    \item generic-refusal and malformed-output rates;
    \item downstream utility and unrelated factual accuracy.
\end{itemize}

\section{Expected Failure Modes}

\paragraph{Weak semantic verifier.}
The guarantee only covers knowledge detected by $C_F$. Improve recall using paraphrase augmentation, entailment models, and adversarial prompt generation.

\paragraph{No safe alternative in the base support.}
If the original model assigns almost all probability to leaking outputs, the tilted target may have too little natural safe mass. Add benign counterfactual responses or a small auxiliary fluency model to expand the candidate set, while retaining the base KL constraint.

\paragraph{High rollout cost.}
Use small candidate sets, few rollouts, cached value estimates, and a learned $h_\psi(s,v)$ critic trained from rollout labels.

\paragraph{Trajectory mismatch.}
Generate training states using the same unmasking/remasking procedure used at inference.

\paragraph{Parallel-token dependency.}
Use one-token microsteps during safe-target construction or keep parallel blocks small.

\section{Summary}

TIDE-U treats MDLM unlearning as minimum-change probability redistribution. It first defines the ideal trajectory distribution as the original model exponentially downweighted by semantic leakage. It then converts that terminal objective into token-level safe distributions through rollout-estimated continuation desirability. The forward-KL training loss decreases logits of tokens with high downstream leakage risk and increases logits of safe alternatives. Separate preservation losses keep ordinary generation close to the original model. Because all token targets derive from the same terminal semantic objective and are trained on inference-like denoising states, the method aims to prevent later denoising steps from reconstructing the forgotten knowledge while avoiding gibberish and hard-coded refusals.

\end{document}
